\subsection{An exact extension of the first norm seed in the two-prime tower}

\begin{statement}{A.156}{Proved}
{An explicit integral norm-two generator at the first changed level}
Put \(d=649=11\cdot59\), and let
\[
 \begin{gathered}
 K_1=\Q(\alpha),\quad\alpha^3=3,\qquad
 K_2=\Q(t),\quad t^9=3d^6=224176165072750803,\\
 \alpha=t^3/d^2,\qquad
 \rho_1=d^8\alpha,\qquad\rho_2=d^2t.
 \end{gathered}
\tag{A.156.1}
\]
These are the actual first two fields in the tower for
\(a=3d^{24}\), since \(\rho_1^3=a\) and \(\rho_2^3=\rho_1\).
The element \(\gamma\) specified below is integral and satisfies
\[
                  N_{K_2/K_1}(\gamma)=\alpha-1,
                 \qquad N_{K_2/\Q}(\gamma)=2.
\tag{A.156.2}
\]
Its principal ideal is exactly the unique prime \(P_{2,0,2}\)
of residue degree one above two.
Consequently
\[
 \beta_0=-2,\qquad\beta_1=1-\alpha,\qquad\beta_2=-\gamma
\tag{A.156.3}
\]
is one actual two-step norm family:
\(N_{K_2/K_1}\beta_2=\beta_1\) and
\(N_{K_1/\Q}\beta_1=\beta_0\).
Each element is a unit away from the primes above two.
In particular the exact class \(e_2\) of A.155 is zero.
The theorem does not claim extension of this particular \(\beta_2\)
to every later level.
\end{statement}

Here is the complete element, using integer blocks to keep its
coefficients legible.  Put \(B=10^{40}\), \(D=d^5=115139273278249\), and
\[
 \gamma=D^{-1}\sum_{i=0}^{8}c_i t^i,\qquad
 c_i=\sigma_i(c_{i,0}+B c_{i,1}+B^2c_{i,2}+B^3c_{i,3}).
\tag{A.156.4}
\]
All entries in the following two tables are nonnegative integers;
\(\sigma_i\) is the displayed sign.
{\scriptsize
\[
\begin{array}{c|c|r|r}
i&\sigma_i&c_{i,0}&c_{i,1}\\\hline
0&+&3254985312134139533616774467688262630288&8865055417849618065965887985680249352793\\
1&-&9818475375430431644865908227448982398242&8608179747455614627828415041243509122607\\
2&+&3605122801677005246908426921985192901559&536374756279981396297437002032616068562\\
3&-&6582044247714884917639856069359167602035&268083660844015591939609172256891704058\\
4&+&1848086138994798713466762895456861371272&9527929021497004532506280934329394304852\\
5&-&1660155711905797346630158742667732636353&5858096507109056107506468482007766062391\\
6&+&1364224149024941749867997586604876840643&3334866010958958032890409384031274427903\\
7&-&6738071505907423414697928082538143249987&502890463303690888204144965696580499069\\
8&+&9079121442640445920992971378334707394599&870805824284542039045774328736682586786\\
\end{array}
\tag{A.156.5}
\]
}
{\scriptsize
\[
\begin{array}{c|r|r}
i&c_{i,2}&c_{i,3}\\\hline
0&5343720612525973055673036975498644633006&2204070986328507041681565181188683\\
1&2270859350602210640019422408003782433074&60035007343402306506629650615224\\
2&5145082889034492655445329518948585338461&1024940204067634301560439951426\\
3&2883610702358533064905204119169277627944&14374364418281344075506156595\\
4&5335812719188460282336095472541049293772&176084999035500536600582192\\
5&3668806510223057878109108366936232864673&1903442969335517751844517\\
6&730091367614578208341668493740717011171&17689767872313241961400\\
7&4998712363626892482906821737836909907510&127178804652686053211\\
8&6920418343839731269190371527301545665855&355955252391989527\\
\end{array}
\tag{A.156.6}
\]
}

\begin{proof}
The defining degree-nine polynomial is Eisenstein at three, so
the fields and degrees in (A.156.1) are as stated.
The equations for the old and new roots follow by taking the
indicated powers; in particular \([K_2:K_1]=3\).

There are two independent exact arithmetic checks of the element.
For the relative norm, put
\[
 A_j=D^{-1}\sum_{k=0}^{2}c_{j+3k}d^{2k}\alpha^k\quad(0\leq j<3),
 \qquad w=d^2\alpha.
\tag{A.156.7}
\]
Thus \(\gamma=A_0+A_1t+A_2t^2\) and \(t^3=w\).
Multiplying the three conjugates gives the literal cubic norm formula
\[
 A_0^3+wA_1^3+w^2A_2^3-3wA_0A_1A_2=\alpha-1.
\tag{A.156.8}
\]
Substitution of (A.156.5)--(A.156.6) and reduction by
\(\alpha^3=3\) verifies all three coefficients in this equality.
Since \(N_{K_1/\Q}(\alpha-1)=2\), it also gives the absolute norm.

To check integrality without using a class-group computation, form
the integer multiplication matrix \(H\), indexed by \(0\leq i,j<9\),
\[
 H_{ij}=\begin{cases}
             c_{i-j},&i\geq j,\\
             (3d^6)c_{9+i-j},&i<j.
          \end{cases}
\tag{A.156.9}
\]
Multiplication by \(\gamma\) in the basis \(1,t,\ldots,t^8\)
is exactly \(D^{-1}H\).
Writing
\(\det(ZI-H)=\sum_{j=0}^{9}h_jZ^{9-j}\), the retained exact
integer calculation verifies
\[
       D^j\mid h_j\ (0\leq j\leq9),\qquad
       \det(H)=2D^9.
\tag{A.156.10}
\]
Thus \(D^{-9}\det(DZI-H)\) is a monic polynomial in \(\Z[Z]\)
with constant coefficient \(-2\), annihilating \(\gamma\).
This proves integrality, and independently gives its norm two.
The full characteristic polynomial, the integer coefficients
\(h_j\), and the literal zero-remainder check are retained in
\path{code/two_inert_norm2_generator_certificate_data.json}.
The program
\path{code/two_inert_norm2_generator_certificate.py}
performs these divisions and determinant calculations over the
rationals and integers, without an analytic or class-group input.

The polynomial defining \(K_2\) has odd discriminant and reduces
modulo two to
\[
 X^9-1=(X-1)(X^2+X+1)(X^6+X^3+1).
\tag{A.156.11}
\]
The last factor is irreducible because the order of two modulo nine
is six, and the middle factor is irreducible of degree two.
Thus there is exactly one prime of norm two, namely the branch
used for \(P_{2,0,2}\) in A.155.
The integral principal ideal \((\gamma)\) has norm two by
(A.156.2), so it is exactly that prime.
Negation has relative norm minus one in the cubic extension.
Equation (A.156.8) therefore proves both norm compatibilities in
(A.156.3), with the fixed first element \(1-\alpha\).
\end{proof}

The element was located by a bounded search for principal relations.
Only its displayed coefficients and the subsequent integer/rational
certificate are used above.  The search's proposed class-group
presentation is not used as a class-number or completeness assertion.
The full search input, request, response and output are retained under
\path{code/two_inert_balanced_level2_magma_v3}.

\paragraph{The specified second element at the next local norm test.}
The successful two-step family in A.156 also gives a literal
example of a representative that cannot be used unchanged at the
following step.  Let \(K_3=\Q(\rho_3)\) with
\(\rho_3^3=\rho_2\), in the same fixed tower.
Then the particular \(\beta_2=-\gamma\) in (A.156.4) satisfies
\[
                 \beta_2\notin N_{K_3/K_2}(K_3^\times).
\tag{A.156.12}
\]
This assertion concerns the specified representative; it does not
exclude another norm preimage of \(\beta_1\).

Here is the full local computation.  At each \(q=11,59\),
the cubic old field \(K_1\) has its degree-one factor and its
unramified quadratic factor.  On the latter the residue of \(\alpha\)
has the indicated equation:
\[
 \begin{array}{c|c|c}
 q&\text{degree-one root of }X^3-3&
          \text{quadratic residue equation for }\alpha\\\hline
 11&9&\alpha^2-2\alpha+4=0\\
 59&49&\alpha^2-10\alpha+41=0.
 \end{array}
\tag{A.156.13}
\]
Both \(K_2/K_1\) and \(K_3/K_2\) are totally ramified cubic
extensions at these branches.  Indeed \(v_{K_1,q}(\rho_1)=8\)
and \(v_{K_2,q}(\rho_2)=8\), prime to three.
The residue fields stay \(\F_{q^2}\), whose multiplicative groups
have three-primary part of order exactly three.

The coefficients \(A_0,A_1\) in (A.156.7) are in
\(\Z[\alpha]\), and \(dA_2\) is in that ring.
Since \(v_{K_2,q}(t)=2\) and \(v_{K_2,q}(q)=3\),
the terms \(A_1t\) and \(A_2t^2\) have positive valuation.
Thus the actual residue of \(\beta_2\) is \(-A_0(\alpha)\).
Substitution of the displayed integer coefficients gives
\[
 \begin{array}{c|c|c|c}
 q&\overline\beta_2&
       \overline\beta_2^{\,(q^2-1)/3}&
       (1-\overline\alpha)^{(q^2-1)/3}\\\hline
 11&6+3\overline\alpha&10+6\overline\alpha&1\\
 59&23+57\overline\alpha&58+6\overline\alpha&1.
 \end{array}
\tag{A.156.14}
\]
The two middle character values are nonidentity cube roots of one;
also \(\overline\beta_2^3=1-\overline\alpha\) in each field,
as required by (A.156.2).
In a totally ramified cubic extension, the residue of the norm
of a unit is its cube.  A norm preimage of the unit \(\beta_2\)
would itself have valuation zero, since the relative residue degree
is one.  The middle column therefore proves that \(\beta_2\)
fails the local norm condition at both of these places, proving
(A.156.12).

For comparison, \(\beta_1=1-\alpha\) meets the unit-norm condition
at these same quadratic places through every higher depth of the
fixed tower.  In a totally ramified extension of degree \(3^j\),
unit norm residues are exactly the \(3^j\)-th powers.
Their image in \(\F_{q^2}^\times\) is the cube subgroup for every
\(j\geq1\), because \(v_3(q^2-1)=1\).
Principal units have all three-power roots in residue characteristic
\(q\ne3\), so a residue norm lifts to an exact unit norm.
The last column of (A.156.14) proves this test for \(\beta_1\).
At the degree-one branches every unit is a three-power scalar norm.
At three use A.140.  At two the distinguished old factor is
\(\Q_2\), with its compatible split root branch.  On the other,
unramified quadratic factor, \(\beta_1\) is a unit and all later
root extensions are unramified, so their unit norms are onto.
At the remaining places use unramified unit norms, and at infinity the
real branch.  Thus the local tests for \(\beta_1\) at all depths
and the failure of this particular \(\beta_2\) are stated on the
same field tower with its actual transition maps.

All residues, cubic characters, coefficient denominators and valuation
bounds in this calculation are independently retained in
\path{code/two_inert_level2_local_norm_certificate.py},
its exact data and its output.

